\begin{proof}[Proof of \cref{obj:lem:calibrated-hybrid-mean-bias}]
For \(x\in[d]\) and \(a\in\{0,1\}\), write
\[
p_x=P\{X=x\},\qquad
s_{ax}=P\{X=x,A=a\},\qquad
q_{ax}=P\{X=x,A=a,Y=1\},
\]
and define the totalized outcome regression by the paper's total real-division convention
\[
\mu_{ax}:=\frac{q_{ax}}{s_{ax}},
\]
so that \(\mu_{ax}=0\) when \(s_{ax}=0\), because real division is totalized with \(r/0=0\).  Put
\[
\vartheta_x:=p_x(\mu_{1x}-\mu_{0x}).
\]
Then \(\tau(P)=\sum_x\vartheta_x\), \(\sum_xp_x=1\), and
\[
|\vartheta_x|
\le p_x,
\]
because \(p_x\ge0\) and the two binary outcome regressions lie in \([0,1]\).  Write \(b_n\) for the binary-length calibration quantity appearing in \(\mathsf C_\epsilon(n,m)\).  Throughout this proof, \(\log\) denotes the natural logarithm; since \(b_n=1+\lceil\log_2 n\rceil\), \(n\ge1\) implies \(\log n=(\log 2)\log_2 n\le b_n\), and hence \(\exp(-12b_n)\le n^{-12}\).  We record the positivity consequences of the finite inequalities defining \(\mathsf C_\epsilon(n,m)\) in \cref{obj:def:mixed-estimator-handle}.  Since \(24\le n\), we have \(n>0\) and \(n+m>0\), so the denominators in the block-proportion bounds are positive.  Those bounds give
\[
\frac1{32}\le \frac un,
\qquad
\frac1{32}\le \frac{t_p}{n+m},
\qquad
\frac1{32}\le \frac t{n+m},
\]
and hence \(u>0\), \(t_p>0\), and \(t>0\).  The scale inequality \(8\le t_pB\), together with \(t_p>0\), then gives \(B>0\).  The calibration-scale definition in \cref{obj:synth_7} makes \(L=L(n)\) a maximum whose first entry is \(2\); in particular \(L\ge2\).  The model condition in \cref{obj:def:model-class} places \(P\) in the occupied-cell overlap class, and \cref{obj:ass:overlap} gives, whenever \(p_x>0\),
\[
\epsilon p_x\le s_{0x},
\qquad
\epsilon p_x\le s_{1x}.
\]
When \(p_x=0\), the identity \(s_{0x}+s_{1x}=p_x\) and nonnegativity of the
two cell masses imply \(s_{0x}=s_{1x}=0\), so the same inequalities remain
valid. % lean: calibrated_hybrid_mean_bias_le

For each cell, let \(J_{0x}^{\mathrm p}\) and \(J_{1x}^{\mathrm p}\) be the two independent pilot arm counts, with
\[
J_{0x}^{\mathrm p}\sim\operatorname{Poi}(t_p s_{0x}),
\qquad
J_{1x}^{\mathrm p}\sim\operatorname{Poi}(t_p s_{1x}),
\qquad
J_x:=J_{0x}^{\mathrm p}+J_{1x}^{\mathrm p}.
\]
Define
\[
\pi_x
:=
\Pr\{J_x\le k_0\},
\qquad
\eta_x
:=
\Pr\{J_x>k_0\}.
\]
Thus \(0\le\pi_x,\eta_x\le1\) and \(\eta_x=1-\pi_x\).  Let \(\mathbb Q_x\) be the independent product law of four armwise Poisson counts
\[
S_{ax}\sim\operatorname{Poi}(u q_{ax}),
\qquad
K_{ax}\sim\operatorname{Poi}(t s_{ax}),
\qquad a\in\{0,1\},
\]
with all four variables mutually independent.  These are ordinary nonnegative Poisson intensities, because \(u,t>0\) and the cell masses \(q_{ax}\) and \(s_{ax}\) are nonnegative.  Define the light and heavy branch means
\[
\Lambda_x
:=
\mathbb E_{\mathbb Q_x}\!\left[
\frac{S_{1x}}u\,\widehat W_{1,L,x}
-
\frac{S_{0x}}u\,\widehat W_{0,L,x}
\right],
\]
and
\[
\Gamma_x
:=
\mathbb E_{\mathbb Q_x}\!\left[
\frac{S_{1x}}u\,\frac{K_{0x}+K_{1x}+1}{K_{1x}+1}
-
\frac{S_{0x}}u\,\frac{K_{0x}+K_{1x}+1}{K_{0x}+1}
\right].
\]
Let \(M_x\) denote the selected cell mean under the independent Poisson count law described in the lemma: the pilot pool selects the light branch on \(\{J_x\le k_0\}\) and the heavy branch on \(\{J_x>k_0\}\), while the outcome and factorial pools are integrated under \(\mathbb Q_x\).  Thus
\[
M_x
=
\mathbb E\!\left[
\mathbf 1_{\{J_x\le k_0\}}\left(
\frac{S_{1x}}u\,\widehat W_{1,L,x}
-
\frac{S_{0x}}u\,\widehat W_{0,L,x}\right)
+
\mathbf 1_{\{J_x>k_0\}}\left(
\frac{S_{1x}}u\,\frac{K_{0x}+K_{1x}+1}{K_{1x}+1}
-
\frac{S_{0x}}u\,\frac{K_{0x}+K_{1x}+1}{K_{0x}+1}\right)
\right].
\]
Independence of the pilot pool from the two evaluation pools gives the exact identity
\[
M_x=\pi_x\Lambda_x+\eta_x\Gamma_x .
\]
% lean: hybridSelectedCellMean_eq_mixture
We now derive the first-moment identity from the armwise Poisson factorial moments, using the factorial-lift definition in \cref{obj:synth_4}.  For \(a\in\{0,1\}\), set
\[
\Lambda_{a,x}
:=
\mathbb E_{\mathbb Q_x}\!\left[
\frac{S_{ax}}u\,\widehat W_{a,L,x}
\right],
\qquad
\Lambda_x=\Lambda_{1,x}-\Lambda_{0,x}.
\]
Under \(\mathbb Q_x\), the outcome count \(S_{ax}\) is independent of the factorial-pool counts, and
\[
\mathbb E_{\mathbb Q_x}S_{ax}=u q_{ax}.
\]
Write the light polynomial as in \cref{obj:synth_4}:
\[
W_{a,L}(r_0,r_1)=\sum_{\rho_0,\rho_1}
 c^{(a)}_{\rho_0,\rho_1}r_0^{\rho_0}r_1^{\rho_1},
\]
with only finitely many nonzero coefficients.  Its factorial lift is
\[
\widehat W_{a,L,x}(K_{0x},K_{1x})
=
\sum_{\rho_0,\rho_1}c^{(a)}_{\rho_0,\rho_1}
\frac{(K_{0x})_{\rho_0}(K_{1x})_{\rho_1}}
     {t^{\rho_0+\rho_1}},
\]
where \((k)_\rho\) is the falling factorial.  If \(K\sim\operatorname{Poi}(ts)\) with \(s\ge0\), then, for every \(\rho\ge0\),
\[
\mathbb E (K)_\rho
=
\sum_{\ell=\rho}^\infty
\frac{\ell!}{(\ell-\rho)!}
 e^{-ts}\frac{(ts)^\ell}{\ell!}
=(ts)^\rho
\sum_{j=0}^\infty e^{-ts}\frac{(ts)^j}{j!}
=(ts)^\rho.
\]
Since \(t>0\) and the two factorial-pool counts are independent Poisson variables with means \(t s_{0x}\) and \(t s_{1x}\), this gives
\[
\mathbb E_{\mathbb Q_x}
\frac{(K_{0x})_{\rho_0}(K_{1x})_{\rho_1}}
     {t^{\rho_0+\rho_1}}
=s_{0x}^{\rho_0}s_{1x}^{\rho_1}.
\]
Termwise summation over the finite coefficient expansion yields the first-moment identity
\[
\mathbb E_{\mathbb Q_x}\widehat W_{a,L,x}=W_{a,L}(s_{0x},s_{1x}),
\]
and therefore
\[
\Lambda_{a,x}
=q_{ax}\,W_{a,L}(s_{0x},s_{1x}).
\]
Using \(s_{0x}+s_{1x}=p_x\) and
\[
W_{a,L}(s_{0x},s_{1x})
=
\frac{p_x}{B}G_L\!\left(\frac{s_{ax}}B\right),
\]
we obtain
\[
\Lambda_x
=
q_{1x}\frac{p_x}{B}G_L\!\left(\frac{s_{1x}}B\right)
-
q_{0x}\frac{p_x}{B}G_L\!\left(\frac{s_{0x}}B\right).
\]
% lean: integral_hybridLightPools

Consider first a light cell, \(p_x\le B\).  If \(p_x=0\), then \(s_{0x}=s_{1x}=q_{0x}=q_{1x}=0\), so \(\Lambda_x=\vartheta_x=0\) and the light-branch bias bound is immediate.  It remains in this paragraph to handle the case \(p_x>0\).  For \(a\in\{0,1\}\), set
\[
z_{ax}:=\frac{s_{ax}}B,
\qquad
E_{ax}:=E_L(z_{ax}).
\]
The overlap inequalities give \(s_{ax}\ge\epsilon p_x>0\) for \(a=0,1\), so \(z_{ax}>0\); the light-cell condition \(p_x\le B\), together with \(s_{ax}\le p_x\) and \(B>0\), gives \(z_{ax}\le1\).  Hence \(z_{ax}\in(0,1]\).  The factorial-moment calculation above has already identified the light-pool first moment.  Since \(L\ge2\) and \(z_{ax}\in(0,1]\), the interval part of \cref{obj:lem:explicit-chebyshev-calibration} applies and gives
\[
z_{ax}G_L(z_{ax})+E_{ax}=1,
\qquad
0\le E_{ax},
\qquad
E_{ax}\le\frac{1}{L^2z_{ax}} .
\]  Since
\(q_{ax}=s_{ax}\mu_{ax}\), the displayed formula for \(\Lambda_x\) becomes, in the case \(p_x>0\),
\[
\Lambda_x
=
p_x\mu_{1x}(1-E_{1x})
-
p_x\mu_{0x}(1-E_{0x}),
\]
and therefore
\[
|\Lambda_x-\vartheta_x|
\le
p_x\mu_{1x}E_{1x}+p_x\mu_{0x}E_{0x}
\le
\frac{p_xB}{L^2s_{1x}}+\frac{p_xB}{L^2s_{0x}}
\le
\frac{2B}{\epsilon L^2}.
\]
Together with the already settled case \(p_x=0\), this proves the light-branch bound for every light cell. % lean: light_exact_mean_bias_le

For the heavy branch, decompose the branch mean armwise by setting
\[
\Gamma_{a,x}
:=
\mathbb E_{\mathbb Q_x}\!\left[
\frac{S_{ax}}u\left(1+\frac{K_{1-a,x}}{K_{ax}+1}\right)
\right],
\qquad a\in\{0,1\},
\]
so that \(\Gamma_x=\Gamma_{1,x}-\Gamma_{0,x}\).  By independence of \(S_{ax}\) from the factorial-pool counts and \(\mathbb E_{\mathbb Q_x}S_{ax}=uq_{ax}\),
\[
\Gamma_{a,x}
=
q_{ax}\,
\mathbb E\!\left[1+\frac{K_{1-a,x}}{K_{ax}+1}\right].
\]
Substituting the cell quantities \(p_x=s_{0x}+s_{1x}\), \(s_{ax}\), \(q_{ax}\), and \(\mu_{ax}=q_{ax}/s_{ax}\) in \cref{obj:lem:poisson-inverse-arm-mean-formula}, and using the total real-division simplification \(u/u=1\) that follows from \(u>0\), gives
\[
q_{ax}\,
\mathbb E\!\left[1+\frac{K_{1-a,x}}{K_{ax}+1}\right]
=
p_x\mu_{ax}-(p_x-s_{ax})\mu_{ax}\exp(-ts_{ax}),
\]
whose left side is \(\Gamma_{a,x}\).
The totalized zero-arm cases are covered as follows: if \(s_{ax}=0\), then \(0\le q_{ax}\le s_{ax}\) gives \(q_{ax}=0\), so \(\Gamma_{a,x}=0\), while \(\mu_{ax}=q_{ax}/s_{ax}=0\) and the right side is also zero.  If \(s_{ax}>0\), the quotient is the ordinary outcome regression and the same identity applies directly.  Thus, for each arm,
\[
|\Gamma_{a,x}-p_x\mu_{ax}|
=
(p_x-s_{ax})\mu_{ax}\exp(-ts_{ax}).
\]
If \(p_x=0\), this expression is zero.  If \(p_x>0\), occupied-cell overlap gives \(s_{ax}\ge\epsilon p_x\); together with \(0\le p_x-s_{ax}\le p_x\) and \(0\le\mu_{ax}\le1\), this yields
\[
|\Gamma_{a,x}-p_x\mu_{ax}|
\le
p_x\exp(-\epsilon t p_x).
\]
Summing the two armwise errors gives
\[
|\Gamma_x-\vartheta_x|
\le
2p_x\exp(-\epsilon t p_x).
\]
% lean: integral_hybridHeavyPools_bias_le

We next derive the pilot-selection estimate used on light cells. Since the two pilot arm counts are independent Poisson variables with means
\(t_p s_{0x}\) and \(t_p s_{1x}\), their sum satisfies
\[
J_x=J_{0x}^{\mathrm p}+J_{1x}^{\mathrm p}
\sim
\operatorname{Poi}\bigl(t_p(s_{0x}+s_{1x})\bigr)
=
\operatorname{Poi}(t_pp_x).
\]
Let
\[
\lambda_x^{\mathrm p}:=t_p p_x.
\]
For any
\(\alpha_{\mathrm{tail}}>0\), the exponential Markov bound applied to
\((1+\alpha_{\mathrm{tail}})^{J_x}\) gives
\[
\begin{aligned}
\Pr\{J_x>k_0\}\exp(-\alpha_{\mathrm{tail}}\lambda_x^{\mathrm p})
&\le
\frac{\mathbb E(1+\alpha_{\mathrm{tail}})^{J_x}}
     {(1+\alpha_{\mathrm{tail}})^{k_0}}
\exp(-\alpha_{\mathrm{tail}}\lambda_x^{\mathrm p}) \\
&=
\left(1+\alpha_{\mathrm{tail}}\right)^{-k_0},
\end{aligned}
\]
because the Poisson probability generating function gives
\(\mathbb E(1+\alpha_{\mathrm{tail}})^{J_x}
=\exp(\alpha_{\mathrm{tail}}\lambda_x^{\mathrm p})\).  Taking
\(\alpha_{\mathrm{tail}}=\epsilon t/t_p\), which is positive since the
calibration predicate gives \(t_p>0\) and \(t>0\), yields
\[
\eta_x\exp(-\epsilon t p_x)
\le
\left(1+\epsilon\frac{t}{t_p}\right)^{-k_0}.
\]
By \(\bar\epsilon\le\epsilon\) from \cref{obj:synth_7} and \(t/t_p\ge0\),
\[
\left(1+\bar\epsilon\frac{t}{t_p}\right)^{k_0}
\le
\left(1+\epsilon\frac{t}{t_p}\right)^{k_0}.
\]
The calibration inequality in \cref{obj:def:mixed-estimator-handle} gives
\[
n^{12}\le\left(1+\bar\epsilon\frac{t}{t_p}\right)^{k_0},
\]
and hence, with positive denominators,
\[
\left(1+\epsilon\frac{t}{t_p}\right)^{-k_0}
\le
n^{-12}.
\]
Combining these estimates with the mixture identity yields, for \(p_x\le B\),
\[
|M_x-\vartheta_x|
\le
\frac{2B}{\epsilon L^2}+2p_x n^{-12}.
\]
% lean: calibrated_light_cell_heavy_probability_mul_exp_le

Now consider a heavy cell, \(B<p_x\).  Put
\[
\lambda_x^{\mathrm p}:=t_p p_x.
\]
We derive the light-pilot lower-tail estimate directly.  By the definition above, \(J_x=J_{0x}^{\mathrm p}+J_{1x}^{\mathrm p}\), and independence of the two pilot arm counts gives
\[
J_x\sim\operatorname{Poi}\bigl(t_p(s_{0x}+s_{1x})\bigr)
=
\operatorname{Poi}(\lambda_x^{\mathrm p}).
\]
Moreover, \cref{obj:def:mixed-estimator-handle} defines the pilot threshold by
\[
k_0=\left\lfloor\frac{t_pB}{4}\right\rfloor.
\]
Because \(t_pB/4\ge0\), the defining floor inequality, after casting the
natural-valued threshold \(k_0\) to the reals, gives \(k_0\le t_pB/4\).
The heavy-cell inequality \(B<p_x\), multiplied by the positive number
\(t_p/4\), gives \(t_pB/4< t_pp_x/4=\lambda_x^{\mathrm p}/4\).  Thus
\[
k_0\le \frac{t_pB}{4}<\frac{\lambda_x^{\mathrm p}}4.
\]
With \(\theta=\log 2\), apply exponential Markov directly to
\(e^{-\theta J_x}\).  Since the event \(J_x\le k_0\) is the same as
\(e^{-\theta J_x}\ge e^{-\theta k_0}\), with the natural number \(k_0\)
cast to a real in the exponential, the floor inequality above gives
\[
\begin{aligned}
\pi_x
&=\Pr\{J_x\le k_0\} \\
&\le
\exp\!\left(\theta k_0
+\lambda_x^{\mathrm p}(e^{-\theta}-1)\right) \\
&\le
\exp\!\left(\frac{\theta\lambda_x^{\mathrm p}}4
+\lambda_x^{\mathrm p}(e^{-\theta}-1)\right) \\
&=
\exp\!\left(\left(\frac{\theta}{4}-\frac12\right)\lambda_x^{\mathrm p}\right)
\le
\exp\!\left(-\frac{\lambda_x^{\mathrm p}}4\right),
\end{aligned}
\]
since \(e^{-\theta}=1/2\) and \(\theta\le1\).  The tail-calibration inequalities in \cref{obj:def:mixed-estimator-handle} include
\[
15L+12b_n\le \frac{t_pB}{4},
\]
so
\[
12b_n\le 15L+12b_n\le \frac{t_pB}{4}<\frac{\lambda_x^{\mathrm p}}4.
\]
Consequently
\[
\pi_x\le \exp(-12b_n)\le n^{-12},
\]
where the last comparison is the binary-length consequence recorded at the start of the proof. % lean: calibrated_heavy_cell_pilot_le

We next prove the auxiliary estimate for the wrong-light contribution on the same heavy cell:
\[
\pi_x|\Lambda_x|\le p_x n^{-12}.
\]
Put
\[
y_x:=\frac{p_x}{B},
\]
so \(y_x\ge1\).  First we bound the absolute light-branch mean.  The first-moment identity gives
\[
\Lambda_x
=
q_{1x}\frac{p_x}{B}G_L\!\left(\frac{s_{1x}}B\right)
-
q_{0x}\frac{p_x}{B}G_L\!\left(\frac{s_{0x}}B\right).
\]
For each arm, \(0\le s_{ax}\le p_x\) and \(0\le q_{ax}\le p_x\).  If
\[
G_L(z)=\sum_{\rho=0}^{L-2} g_\rho z^\rho,
\]
then the degree bound in \cref{obj:lem:chebyshev-factorial-certificate} justifies the displayed finite expansion with no terms beyond degree \(L-2\), while \cref{obj:lem:explicit-chebyshev-calibration} gives
\(2\|G_L\|_1\le14^L\).  Since \(0\le s_{ax}/B\le y_x\) and \(y_x\ge1\), each displayed exponent satisfies
\[
\left(\frac{s_{ax}}B\right)^\rho\le y_x^\rho\le(1+y_x)^\rho\le(1+y_x)^L.
\]
Multiplying this monomial bound by \(|g_\rho|\) and summing over the displayed finite support gives
\[
\left|G_L\!\left(\frac{s_{ax}}B\right)\right|
\le \sum_{\rho=0}^{L-2}|g_\rho|\left(\frac{s_{ax}}B\right)^\rho
\le \|G_L\|_1(1+y_x)^L.
\]
Also \(y_x\le(1+y_x)^L\), because \(L\ge2\) and \(y_x\ge0\).  Therefore
\[
|\Lambda_x|
\le
2p_xy_x\|G_L\|_1(1+y_x)^L
\le
p_x\,14^L(1+y_x)^{2L}.
\]
This coefficient-envelope argument uses the exact first-moment identity above, together with the degree and coefficient bounds for \(G_L\), and yields the displayed bound without invoking any second-moment estimate.
% lean: abs_integral_hybridLightPools_heavy_le
The lower-tail calculation above, kept in exponential form, gives
\[
\pi_x\le\exp\!\left(-\frac{t_pp_x}{4}\right).
\]
The explicit scale definition in \cref{obj:synth_7} gives
\[
B=\frac{H_\epsilon L}{\min(t_p,t)}.
\]
Since \(t_p>0\) and \(t>0\), their minimum is positive and the scale identity
gives \(H_\epsilon L=B\min(t_p,t)\).  As \(\min(t_p,t)\le t_p\) and
\(B>0\), multiplication preserves the inequality, so \(H_\epsilon L\le
t_pB\).  Thus, with \(p_x=By_x\),
\[
\exp\!\left(-\frac{t_pp_x}{4}\right)
\le
\exp\!\left(-\frac{H_\epsilon Ly_x}{4}\right).
\]
The environment \cref{obj:synth_7} gives \(A_\star=7056\),
\(c^\circ=1/256\), and
\[
\frac{H_\epsilon}{4}\ge 15+\frac{24}{c^\circ}.
\]
Elementary arithmetic applied to the displayed value of \(A_\star\) gives
\(14\le A_\star<e^{13}\).  Using these consequences, \(y_x\ge1\), and
\(\log(1+y_x)\le y_x\), we obtain
\[
L\log A_\star+2L\log(1+y_x)-\frac{H_\epsilon Ly_x}{4}
\le
13Ly_x+2Ly_x-\left(15+\frac{24}{c^\circ}\right)Ly_x
=
-\frac{24Ly_x}{c^\circ}.
\]
Exponentiating gives the calibrated absorption bound
\[
A_\star^L(1+y_x)^{2L}
\exp\!\left(-\frac{H_\epsilon Ly_x}{4}\right)
\le
\exp\!\left(-\frac{24Ly_x}{c^\circ}\right).
\]
Combining the envelope for \(\Lambda_x\), the pilot-tail display, and \(14^L\le A_\star^L\),
\[
\pi_x|\Lambda_x|
\le
p_x\exp\!\left(-\frac{24Ly_x}{c^\circ}\right).
\]
Finally, the degree lower bound \(c^\circ b_n/2\le L\) in \(\mathsf C_\epsilon(n,m)\), listed in \cref{obj:def:mixed-estimator-handle}, together with \(y_x\ge1\), implies
\[
12b_n\le \frac{24Ly_x}{c^\circ}.
\]
The binary-length consequence recorded at the start of the proof gives
\[
\exp\!\left(-\frac{24Ly_x}{c^\circ}\right)
\le
\exp(-12b_n)
\le
n^{-12}.
\]
Combining the last two displays proves
\[
\pi_x|\Lambda_x|\le p_xn^{-12}.
\]
% lean: calibrated_heavy_cell_light_mean_le

The intrinsic heavy-branch exponential is also absorbed on a heavy cell:
\[
\exp(-\epsilon t p_x)
\le
n^{-12}.
\]
Indeed, \(B<p_x\), \(\bar\epsilon\le\epsilon\), and the tail-calibration inequality in \cref{obj:def:mixed-estimator-handle},
\[
12b_n\le \bar\epsilon tB,
\]
give \(12b_n\le \epsilon t p_x\), and hence, by the binary-length consequence recorded at the start of the proof,
\[
\exp(-\epsilon t p_x)\le \exp(-12b_n)\le n^{-12}.
\]
Therefore
\[
\eta_x|\Gamma_x-\vartheta_x|
\le
2p_xn^{-12}.
\]
Together with \(\pi_x|\vartheta_x|\le p_xn^{-12}\), the decomposition
\[
M_x-\vartheta_x
=
\pi_x\Lambda_x-\pi_x\vartheta_x+\eta_x(\Gamma_x-\vartheta_x)
\]
gives, for \(B<p_x\),
\[
|M_x-\vartheta_x|
\le
4p_xn^{-12}.
\]
% lean: calibrated_heavy_cell_bias_exp_le

The alternatives \(p_x\le B\) and \(B<p_x\) are disjoint and exhaustive, so
they partition \([d]\).  Summing the two cellwise bounds over these two regimes,
\[
\begin{aligned}
\left|
\sum_{x\in[d]}M_x-\tau(P)
\right|
&=
\left|
\sum_{x\in[d]}\bigl(M_x-\vartheta_x\bigr)
\right|  \\
&\le
\sum_{x\in[d]}
\left|M_x-\vartheta_x\right| \\
&\le
\sum_{x:\,p_x\le B}
\left(\frac{2B}{\epsilon L^2}+2p_xn^{-12}\right)
+
\sum_{x:\,B<p_x}4p_xn^{-12} \\
&\le
\frac{2dB}{\epsilon L^2}
+4n^{-12}\sum_{x\in[d]}p_x \\
&=
\frac{2dB}{\epsilon L^2}+4n^{-12}
\le
\frac{2dB}{\epsilon L^2}+6n^{-12}.
\end{aligned}
\]
This is the claimed bound. % lean: calibrated_hybrid_mean_bias_le
\end{proof}
